\section{Introduction}
\label{sec:intro}

Dermoscopic classifiers now rival specialists on curated benchmarks~\cite{esteva2017dermatologist}, and
explanation heatmaps are a common way to argue that such a model is trustworthy: if the heatmap falls on
the lesion and the highlighted pixels drive the prediction, the model is said to be looking at the right
thing. Two benchmark families operationalize this argument. Localization scores a heatmap against an
expert lesion mask, by intersection over union (IoU) or by the pointing game~\cite{zhang2016topdown}.
Faithfulness removes or inserts pixels in heatmap order and tracks the predicted-class
probability~\cite{petsiuk2018rise,samek2017evaluating}. Both are routinely used to compare architectures,
often with each architecture paired with its customary explainer: Grad-CAM~\cite{selvaraju2017gradcam}
for CNNs and attention rollout~\cite{abnar2020attention} for Vision Transformers~\cite{kozachok2026melanoscope}.

This paper asks whether these benchmarks, as commonly applied to dermoscopy, measure the model. Two
properties of the setting make the answer non-obvious. First, lesions in HAM10000~\cite{tschandl2018ham10000}
are usually photographed near the image center, so a heatmap that ignores the image may already score
well on localization. Second, attention rollout does not depend on the target class~\cite{chefer2021transformer},
whereas deletion/insertion are computed on the predicted-class probability, which structurally favors a
class-specific explainer such as Grad-CAM. A comparison of ``CNN with Grad-CAM'' against ``ViT with
rollout'' therefore compares two model-explainer systems, not two architectures.

\textbf{Contributions.} (1) A single protocol that scores classification, faithfulness and localization
on the same checkpoints and the same \todo{1,002} test images of a fixed lesion-grouped split, with three
training seeds per model and paired image-level confidence intervals. (2) Two model-free controls, a
random pixel order and a fixed centered heatmap, reported as rows in the main tables. \todo{(3) Finding:
the center prior [matches/exceeds] Grad-CAM on IoU and pointing game; numbers from
results/final/tables/localization.tex.} (4) One shared, class-specific explainer, Integrated
Gradients~\cite{sundararajan2017axiomatic}, applied to both architectures, which \todo{[does/does not]}
preserve the architecture ranking obtained with the customary explainers.
\TODO{Two sentences on FABRIC fit: trustworthy evaluation of biomedical imaging models; why evaluation
artifacts matter before clinical studies. Keep claims at research level, no clinical-readiness claims.}
